\documentclass[10pt,a4paper]{amsart}
\usepackage[margin=19mm]{geometry}
\usepackage{amsmath,amssymb,mathtools}
\usepackage[scr=rsfs,cal=euler]{mathalfa}
\usepackage{xfrac}
\usepackage[unicode,hidelinks,bookmarks=false]{hyperref}
\newcommand{\ie}{i.e.\ }
\newcommand{\cf}{cf.\ }
\newcommand{\resp}{respectively\ }
\newcommand{\Subsection}{\subsection}
\providecommand{\nobreakdash}{\nobreak\hbox{-}}
\setlength{\emergencystretch}{2em}
\multlinegap0pt
\usepackage{amssymb}
\usepackage{mathtools}
\usepackage{float}
\newtheorem{lemma}{Lemma}
\title{Almost-Orthogonality in \texorpdfstring{$L^p$}{Lp} Spaces: A Case Study with Grok}
\author{Ziang Chen, Jaume de Dios Pont, Paata Ivanisvili, Jose Madrid, Haozhu Wang}
\date{Source: arXiv:2605.05192v1 (CC BY 4.0)}
\begin{document}
\maketitle

\noindent\emph{Source and licence.} Almost-Orthogonality in \texorpdfstring{$L^p$}{Lp} Spaces: A Case Study with Grok. Version arXiv:2605.05192v1 (CC BY 4.0), distributed by arXiv under the Creative Commons Attribution 4.0 International licence, http://creativecommons.org/licenses/by/4.0/, as recorded in the arXiv licence metadata for this exact version and shown on the abstract page https://arxiv.org/abs/2605.05192v1. Full licence notice: \texttt{licenses/arxiv-2605.05192-CC-BY-4.0.txt}.\par\smallskip

\noindent\textit{Editorial scope.} Counted unit: the article's lemma lem:Hr-positive (source0.tex lines 2002--2014). The article's complete original proof of it is delivered with it (source0.tex lines 2016--2264, one \texttt{proof} environment of 249 lines). No mathematics is written, repaired or re-derived: the statement and the proof are the article's own lines, in the article's own notation, and no retained line is substituted or re-flowed.  Retained uncounted as the source-local context the proof needs: lines 961--963.  Nothing was omitted from any retained span, so the package carries no omission record.  No retained line contains a citation, so the wrapper carries no bibliography and no bibliography entry.  No retained line contains a figure environment, an image-inclusion command or drawing code, so no image is delivered and the package declares no asset dependency.  Editorial formatting only, in addition: an AMS \texttt{amsart} preamble that restates the article's own declarations and macros used by the retained lines, namely the environments \texttt{lemma} on separate counters, together with the standard packages those lines need.  The retained environments print in this document as Lemma 1 in order of appearance, whereas the article prints the same items with the numbering of its own full document.  The complete native source of the article as distributed in the arXiv e-print for this exact version is bundled unchanged as \texttt{sources/199934/source0.tex}.
\begin{lemma}\label{lem: tecnichal lemma about Q'sA and Hr}
 We have $Q'(s_{A})<0$, and $H(r)>0$.  
\end{lemma}

\begin{lemma}[the remainin part of Lemma~\ref{lem: tecnichal lemma about Q'sA and Hr}: $H(r)>0$]\label{lem:Hr-positive}
Let $p>3$, set
\[
r=\frac{p}{p+1},
\qquad
c=c(p)=\frac{2\ln 2}{(p-2)\ln 3+2\ln 2}.
\]
Let $Q$ be the pseudo--polynomial defined above (so that $P''(s)=s^{p-4}Q(s)$).
Then
\[
r\,Q''(r)-(p+1)\,Q'(r)>0.
\]
\end{lemma}

\begin{proof}
Define
\[
H(s):=sQ''(s)-(p+1)Q'(s).
\]
For a monomial $a s^\alpha$ one has
\[
\frac{d}{ds}(a s^\alpha)=a\alpha s^{\alpha-1},\qquad
\frac{d^2}{ds^2}(a s^\alpha)=a\alpha(\alpha-1)s^{\alpha-2},
\]
hence
\[
s\frac{d^2}{ds^2}(a s^\alpha)-(p+1)\frac{d}{ds}(a s^\alpha)
= a\alpha(\alpha-p-2)s^{\alpha-1}.
\]
Applying this term-by-term to the explicit formula for $Q$ and then substituting
$r=\frac{p}{p+1}$, we obtain the exact identity
\begin{equation}\label{eq:H-reduction}
H(r)=\frac{2(p-1)}{(p+1)^3}\,S(p,c,u),
\qquad
u:=r^{p-1}=\Bigl(\frac{p}{p+1}\Bigr)^{p-1}\in(0,1),
\end{equation}
where $S$ is a quadratic polynomial in $u$:
\[
S(p,c,u)=A(p,c)+B(p,c)\,u+C(p,c)\,u^2,
\]
with
\begin{align*}
A(p,c)={}&
30cp^7+30cp^6-54cp^5-70cp^4+8cp^3+40cp^2+16cp\\
&\quad +9p^7-24p^6-66p^5+57p^3+16p^2-16p-8,\\[2mm]
B(p,c)={}&
-72cp^7-64cp^6+157cp^5+148cp^4-80cp^3-77cp^2\\
&\quad -18p^7+45p^6+156p^5+8p^4-148p^3-51p^2+16p,\\[2mm]
C(p,c)={}&
-108cp^8+138cp^7+217cp^6-307cp^5-154cp^4+199cp^3+39cp^2-36cp\\
&\quad +135p^7-33p^6-294p^5+32p^4+207p^3+13p^2-36p.
\end{align*}
Since $\frac{2(p-1)}{(p+1)^3}>0$ for $p>3$, it suffices to prove
\begin{equation}\label{eq:S-positive-goal}
S(p,c(p),u(p))>0.
\end{equation}

%------------------- Step 1: rational bounds for c(p) -------------------------
\medskip
\noindent\textbf{Step 1: rational bounds for $c(p)$.}
Write $a:=\log_2 3=\frac{\ln 3}{\ln 2}$. Since
\[
3^5=243<256=2^8
\quad\text{and}\quad
3^{12}=531441>524288=2^{19},
\]
we have
\[
\frac{19}{12}<a<\frac85.
\]
Using $c(p)=\dfrac{2}{\,2+(p-2)a\,}$ and that $c$ is decreasing in $a$, it follows that
\begin{equation}\label{eq:c-bounds-sharp}
c_-:=\frac{5}{4p-3}<c(p)<c_+:=\frac{24}{19p-14}.
\end{equation}

%------------------- Step 2: an explicit upper bound for u(p) -----------------
\medskip
\noindent\textbf{Step 2: an explicit upper bound for $u(p)$.}
Let $n:=p-1>2$ and $x:=1/p\ge 0$. Consider
\[
f(x):=(1+x)^n-\Bigl(1+nx+\frac{n(n-1)}{2}x^2\Bigr).
\]
Then $f(0)=f'(0)=f''(0)=0$, and for $x\ge0$,
\[
f'''(x)=n(n-1)(n-2)(1+x)^{n-3}\ge 0.
\]
Hence $f''$ is increasing with $f''(0)=0$, so $f''\ge0$; thus $f'$ is increasing
with $f'(0)=0$, so $f'\ge0$; thus $f$ is increasing with $f(0)=0$, so $f\ge0$.
Therefore, for $x\ge0$,
\[
(1+x)^n\ge 1+nx+\frac{n(n-1)}{2}x^2.
\]
With $x=1/p$ and $n=p-1$ this gives
\[
\Bigl(1+\frac1p\Bigr)^{p-1}
\ge
1+\frac{p-1}{p}+\frac{(p-1)(p-2)}{2p^2}
=
\frac{5p^2-5p+2}{2p^2}.
\]
Taking reciprocals yields the bound
\begin{equation}\label{eq:u-upper}
u=\Bigl(\frac{p}{p+1}\Bigr)^{p-1}
=\frac{1}{(1+1/p)^{p-1}}
\le
U(p):=\frac{2p^2}{5p^2-5p+2}.
\end{equation}

%------------------- Step 3: S decreases in u on [0,U(p)] ---------------------
\medskip
\noindent\textbf{Step 3: $S$ is decreasing in $u$ on $[0,U(p)]$ for $c=c(p)$.}
Since $S(p,c,u)=A+B u+C u^2$ is quadratic in $u$,
\[
\frac{\partial S}{\partial u}(p,c,u)=B(p,c)+2C(p,c)\,u
\]
is \emph{linear} in $u$. Hence it suffices to show
\begin{equation}\label{eq:Su-endpoints}
\frac{\partial S}{\partial u}(p,c(p),0)<0
\qquad\text{and}\qquad
\frac{\partial S}{\partial u}(p,c(p),U(p))<0.
\end{equation}

\smallskip
\emph{(i) The endpoint $u=0$.}
We have $\frac{\partial S}{\partial u}(p,c,0)=B(p,c)$. The coefficient of $c$ in $B$
equals
\[
-72p^7-64p^6+157p^5+148p^4-80p^3-77p^2
=
-p^2\Bigl(72p^5+64p^4-157p^3-148p^2+80p+77\Bigr).
\]
For $p=3+x$ with $x\ge0$,
\[
72p^5+64p^4-157p^3-148p^2+80p+77
=
72x^5+1144x^4+7091x^3+21335x^2+31025x+17426>0,
\]
so $B(p,c)$ is strictly decreasing in $c$ for all $p\ge3$.
Using $c(p)\ge c_-$ from \eqref{eq:c-bounds-sharp} gives $B(p,c(p))\le B(p,c_-)$.
A direct substitution $c=c_-=\frac{5}{4p-3}$ yields
\[
B(p,c_-)=
-\frac{p}{4p-3}\,N_B(p),
\]
where
\[
N_B(p):=72p^7+126p^6-169p^5-349p^4-124p^3+160p^2+168p+48.
\]
For $p=3+x$ with $x\ge0$,
\[
N_B(3+x)=72x^7+1638x^6+15707x^5+82166x^4+252638x^3+455074x^2+442767x+178626>0,
\]
hence $B(p,c_-)\!<0$, and therefore $B(p,c(p))<0$.

\smallskip
\emph{(ii) The endpoint $u=U(p)$.}
Set
\[
D(p,c):=B(p,c)+2C(p,c)\,U(p)=\frac{\partial S}{\partial u}(p,c,U(p)).
\]
A simplification gives
\[
D(p,c)=
-\frac{p}{5p^2-5p+2}\Bigl(M(p)\,c + K(p)\Bigr),
\]
with
\[
M(p):=432p^9-192p^8-908p^7+267p^6+789p^5+30p^4-467p^3-81p^2+154p.
\]
For $p=3+x$ with $x\ge0$,
\begin{align*}
\frac{M(3+x)}{3+x}
=&
432x^8+10176x^7+103924x^6+600819x^5+2150214x^4\\
&+4877544x^3
+6849487x^2+5445906x+1877824>0,
\end{align*}
so $M(p)>0$ for $p\ge3$, and thus $D(p,c)$ is strictly decreasing in $c$.
Using again $c(p)\ge c_-$ gives $D(p,c(p))\le D(p,c_-)$.
Substituting $c=c_-=\frac{5}{4p-3}$ yields
\[
D(p,c_-)= -\frac{p}{(4p-3)(5p^2-5p+2)}\,N_D(p),
\]
where
\[
N_D(p):=360p^9-342p^8-1363p^7+1452p^6+939p^5-982p^4-256p^3+8p^2+96p+96.
\]
For $p=3+x$ with $x\ge0$,
\begin{align*}
N_D(3+x)=&
360x^9+9378x^8+107069x^7+703125x^6+2926524x^5+8004428x^4\\&
+14383469x^3+16369613x^2+10703106x+3061824>0,
\end{align*}
hence $D(p,c_-)\!<0$, and therefore $D(p,c(p))<0$.

Combining (i) and (ii) proves \eqref{eq:Su-endpoints}. Since $\partial S/\partial u$
is linear in $u$, it follows that $\partial S/\partial u<0$ for all $u\in[0,U(p)]$,
so $S(p,c(p),u)$ is strictly decreasing on $[0,U(p)]$. Using $u(p)\le U(p)$
from \eqref{eq:u-upper} we obtain
\begin{equation}\label{eq:S-lower-by-U}
S(p,c(p),u(p))\ge S(p,c(p),U(p)).
\end{equation}

%------------------- Step 4: S(p,c,U) decreases in c and replace c by c_+ -----
\medskip
\noindent\textbf{Step 4: replace $c(p)$ by the upper bound $c_+$.}
Since $S$ is affine in $c$, it suffices to compute its $c$--coefficient at $u=U(p)$.
A direct simplification gives
\[
\frac{\partial S}{\partial c}\bigl(p,c,U(p)\bigr)
=
-\frac{2p(p-1)}{(5p^2-5p+2)^2}\,N_2(p),
\]
where
\[
N_2(p):=216p^{10}-75p^9-174p^8+229p^7-140p^6-107p^5+179p^4+54p^3-72p^2-48p+32.
\]
For $p=3+x$ with $x\ge0$,
\begin{align*}
N_2(3+x)=&
216x^{10}+6405x^9+85281x^8+671593x^7+3464881x^6+12239092x^5\\&
+29980589x^4+50292735x^3+55296009x^2+35983743x+10524704>0,
\end{align*}
hence $\frac{\partial S}{\partial c}(p,c,U(p))<0$ for all $p>3$.
Therefore $S(p,c,U(p))$ is strictly decreasing in $c$.
Using $c(p)\le c_+$ from \eqref{eq:c-bounds-sharp} yields
\begin{equation}\label{eq:S-lower-by-cplus}
S(p,c(p),U(p))\ge S(p,c_+,U(p)),
\qquad
c_+=\frac{24}{19p-14}.
\end{equation}

%------------------- Step 5: explicit positivity at (c_+,U(p)) ----------------
\medskip
\noindent\textbf{Step 5: positivity at $(c_+,U(p))$.}
Substituting $c=c_+$ and $u=U(p)$ into $S=A+Bu+Cu^2$ simplifies to
\[
S(p,c_+,U(p))
=
\frac{(p-2)(p-1)}{(19p-14)(5p^2-5p+2)^2}\,R(p),
\]
where
\[
R(p):=747p^{10}-2469p^9+102p^8+4954p^7-1385p^6-2997p^5+2124p^4+380p^3-1096p^2+128p+224.
\]
For $p=3+x$ with $x\ge0$,
\begin{align*}
R(3+x)=&
747x^{10}+19941x^9+235974x^8+1627726x^7+7235131x^6+21607281x^5\\&
+43792452x^4+59291840x^3+51146048x^2+25304972x+5451008>0,
\end{align*}
hence $R(p)>0$ for all $p\ge3$, and therefore $S(p,c_+,U(p))>0$ for all $p>3$.

Finally, combining \eqref{eq:S-lower-by-U} and \eqref{eq:S-lower-by-cplus} gives
\[
S(p,c(p),u(p))\ge S(p,c(p),U(p))\ge S(p,c_+,U(p))>0,
\]
which proves \eqref{eq:S-positive-goal}. Returning to \eqref{eq:H-reduction} yields
$H(r)>0$, i.e.
\[
rQ''(r)-(p+1)Q'(r)>0.
\]
\end{proof}

\end{document}
